\section{The Historical Backtesting Protocol}
\label{sec:protocol}

The protocol evaluates a set of frozen questions $Q = \{q_1,\dots,q_n\}$
against a future corpus. It has six steps; each is fully specified so
that two groups running the same instance obtain the same measurement.

\subsection{Step 1: Choose a cutoff}
A historical date $T$ (Astronomy~v1: 2020-12-31) splits the literature
into a \emph{past corpus} \corpusA{} (everything available up to $T$)
and a \emph{future window} realized as an isolated corpus \corpusB{}
(strictly after $T$). The cutoff must be far enough in the past for the
community to have had time to act---we recommend $\geq 4$ years---and
recent enough that the past corpus reflects a modern research frontier.

\subsection{Step 2: Generate questions}
Any method may generate questions: an evidence-graph pipeline, a
prompted LLM, a heuristic over citation statistics, a human expert. The
only requirements are (i) the generator consumes \emph{only} \corpusA{}
evidence, and (ii) every question records the pre-cutoff evidence it is
grounded in (\texttt{source\_evidence\_ids}). A submission whose source
evidence postdates $T$ is invalid, mechanically
(\texttt{scripts/validate\_cutoff.py}).

\subsection{Step 3: Freeze}
Questions are serialized---identifier, text, cutoff, generating system,
source evidence, system-assigned rank---with \texttt{frozen:\,true}
\emph{before any access to post-cutoff literature}, and are never edited
afterwards. Freezing is the protocol's load-bearing rule: without it,
question text drifts toward what the evaluator has meanwhile learned the
future contains, and the backtest silently becomes a description of the
future rather than a prediction of it \citep{bailey2014backtest}. In the
released implementation frozen files are append-only and guarded by CI;
editing a released question mints a new versioned instance rather than
overwriting the old one.

\subsection{Step 4: Define the future window}
\corpusB{} is collected under its own frozen manifest (query set, date
window, deduplication rules) and stored separately from \corpusA{};
records from \corpusB{} must never enter the generation pipeline. The
Astronomy~v1 window is 2021--2026. Bounding the window matters for
comparability: ``eventually engaged'' is not a fixed target, but
``engaged within $k$ years'' is.

\subsection{Step 5: Retrieve future evidence}
For each frozen question, the question text and every \corpusB{}
document (title + abstract) are embedded
(\texttt{text-embedding-3-small}); the top-$k$ documents by cosine
similarity ($k=8$) become the candidate evidence. Retrieval is
deliberately fixed and deliberately simple: systems are compared on
their \emph{questions}, not their retrievers, and reviewers can inspect
exactly which documents the judge saw because retrieval records are part
of the release.

\subsection{Step 6: Assess outcomes}
\label{sec:protocol:outcomes}

A judge reads the question and its retrieved candidates and assigns two
\emph{independent} labels (Table~\ref{tab:taxonomy}): the fate of the
question and the fate of its premise.

\begin{table}[t]
  \centering
  \small
  \caption{Outcome taxonomy v1.0. The two dimensions are labeled
  independently.}
  \label{tab:taxonomy}
  \begin{tabular}{@{}llp{8.2cm}@{}}
    \toprule
    Dimension & Label & Meaning \\
    \midrule
    \multirow{4}{*}{\texttt{outcome}}
      & \labelfmt{answered} & Future evidence substantially answers the question (including by refuting its premise) \\
      & \labelfmt{partially\_addressed} & Important, directly relevant progress; core question unresolved \\
      & \labelfmt{posed\_but\_open} & The community independently poses essentially the same question without resolving it \\
      & \labelfmt{not\_addressed} & No meaningful follow-up in the future corpus \\
    \midrule
    \multirow{5}{*}{\texttt{premise\_status}}
      & \labelfmt{supported} & Future evidence confirms the underlying premise \\
      & \labelfmt{refuted} & Future evidence falsifies the underlying premise \\
      & \labelfmt{weakened} & Substantial doubt cast without falsification \\
      & \labelfmt{still\_plausible} & The premise was not directly tested after the cutoff \\
      & \labelfmt{not\_applicable} & The question rests on no contestable premise \\
    \bottomrule
  \end{tabular}
\end{table}

Two design decisions deserve emphasis. First, the dimensions are
separated because the single most informative outcome a backtest can
surface---\emph{the community answered this question by refuting its
premise}---is inexpressible in a flat label set: it is simultaneously a
resolution (\labelfmt{answered}) and a falsification
(\labelfmt{refuted}). Our pilot's headline case
(Section~\ref{sec:results:case}) is exactly of this type. Second,
premise refutation is scored as a \emph{success} of the question, not a
failure: a question that provokes the community into overturning one of
its own published conclusions has done the most a question can do.

The judge operates under hard constraints enforced outside the model:
it may cite only bibcodes from the retrieved candidates (violations are
errors, never silently dropped); topical similarity is explicitly
insufficient---the cited paper must bear on the question's actual test
or premise; unknown labels fall back to the most conservative value and
are flagged. Labels then occupy one of two declared tiers.
\emph{Adjudicated} labels have passed human review under written
guidelines (citation validity, the engagement bar, the outcome/premise
split), with the adjudication log released; the headline labels of a
released instance are required to be of this tier (Astronomy~v1's are).
\emph{Judge-only} labels have not, are marked as such wherever
reported, and are the tier at which this paper's large-$n$ comparative
studies run (Sections~\ref{sec:scaled}--\ref{sec:foresight}). The
tier is part of every result's provenance; conflating them is a
protocol violation.

\subsection{Submissions}
A system is evaluated by submitting a directory containing
\texttt{questions.jsonl} (the frozen questions) and
\texttt{metadata.json} (system description, including a mandatory
declaration of any LLM components and their versions, for contamination
auditing). The benchmark pipeline---isolation checks, retrieval,
judging, adjudication where the tier requires it, metrics,
report---is identical for every
submission; the system whose questions we evaluate in
Section~\ref{sec:results} interacts with the benchmark only through
this interface.
